\documentclass[12pt]{article}
\usepackage{ceai}

\begin{document}
\title{Tutorial 04:  ML Production in Action - 01}
\author{}
\date{}

\maketitle

\section*{Engineering Motivation}

In structural design practice, compressive strength is not merely a continuous variable to be predicted. Instead, it informs discrete engineering decisions such as:

\begin{itemize}
    \item Is this mix design acceptable under normal strength specifications?
    \item Does this mix qualify as high-strength concrete (HSC)?
    \item Can this batch be approved for structural placement?
\end{itemize}

In this tutorial, we formulate a supervised binary classification problem grounded in structural engineering standards.

Rather than predicting compressive strength directly, we will classify concrete into:

\begin{align*}
y =
\begin{cases}
1 & \text{Normal Strength Concrete (NSC)} \\
0 & \text{High Strength Concrete (HSC)}
\end{cases}
\end{align*}

based on the standardized 28-day compressive strength.

\vspace{0.5em}

\noindent
The classification threshold is defined as:

\[
f'_c(28) \le 40 \text{ MPa} \quad \Rightarrow \quad \text{NSC}
\]
\[
f'_c(28) > 40 \text{ MPa} \quad \Rightarrow \quad \text{HSC}
\]

\section*{Dataset Description}

We use the UCI Concrete Compressive Strength dataset (available for the class at 
\href{https://umsystem.instructure.com/courses/407601/files/40676265?module_item_id=11763333}{this link}). 
After downloading, unzip the files and use the data in \texttt{Concrete\_Data.xlsx}. Open the spreadsheet, we will see it is structured as: 

\paragraph{Features}
\begin{enumerate}
    \item Cement (kg/m$^3$)
    \item Blast Furnace Slag (kg/m$^3$)
    \item Fly Ash (kg/m$^3$)
    \item Water (kg/m$^3$)
    \item Superplasticizer (kg/m$^3$)
    \item Coarse Aggregate (kg/m$^3$)
    \item Fine Aggregate (kg/m$^3$)
    \item Age (days)
\end{enumerate}

\paragraph{Target variable}
\begin{itemize}
    \item Measured compressive strength (MPa)
\end{itemize}

This data set serves well for formulating a supervised \textit{regression} problem or a time-varying concrete strength forecasting problem, which will be explored later; yet, for the task at hand (i.e., classification), still with significant engineering meaningfulness, we need to conduct some `manual' processing. 

\section*{Strength Standardization}
\subsection*{Tech Spec}
Concrete strength evolves with curing time. Since the dataset includes specimens tested at different ages, we must normalize strengths to the 28-day equivalent.

According to ACI 209.2R-08, the compressive strength development function may be approximated as:

\[
f'_c(t) = f'_c(28) \cdot \frac{t}{a + bt}
\]

where:
\begin{itemize}
    \item $t$ = age (days),
    \item $a, b$ = empirical constants.
\end{itemize}

For typical normal-weight concrete under standard curing:

\[
a = 4.0, \qquad b = 0.85
\]

Rearranging to compute standardized 28-day strength:

\[
f'_c(28) = f'_c(t) \cdot \frac{a + bt}{t}
\]

Thus, for each observation:

\begin{enumerate}
    \item Extract measured strength $f'_c(t)$,
    \item Extract age $t$,
    \item Compute equivalent $f'_c(28)$ using the above equation.
\end{enumerate}

\subsection*{Binary Label Construction}

After computing the standardized 28-day strength:

\[
y =
\begin{cases}
1 & \text{if } f'_c(28) \le 40 \text{ MPa} \\
0 & \text{if } f'_c(28) > 40 \text{ MPa}
\end{cases}
\]

\subsection*{Processing and Visualization}
At this moment, let's conduct some processing and visualization as a Civil Engineer - we wear our Machine Learning Production Engineer hat later. Let complete in Excel:

\begin{itemize}
    \item Calculate the 28-day concrete strength.
    \item Assign binary labels (i.e., `1' - NSC; `0' - HSC). Or, if you like, you can use the labels `NSC' and `HSC' directly.
    \item Select one feature, and plot the histogram of the feature for NSC and HSC in one plot.
    \item Select any two features and plot scatter-plots for NSC and HSC in one plot. 
    \end{itemize}
    
Discussion items:
\begin{itemize}
    \item Do you see \textit{patterns}?
    \item How might changing the threshold affect class balance - going back to the binary scheme - what if we change the threshold \footnote{In modern concrete code, normally 42 MPa is the threshold}? 
    \item Is this threshold universally accepted?
\end{itemize}

\section*{Exploratory Data Analysis in Colab}
Now we will switch to Colab - we will `vibe' coding yet with a design in mind. Before embarking on ML production, let's perform data processing and visualization (as done above), considering the following tasks and keywords for the prompts. 

First, place the updated \texttt{Concrete\_Data} data (with 28-day strength and label) in your Google Drive folder, specify its path, and request that it be mounted. Then practice loading it in (e.g., prompt: load the spreadsheet \texttt{Concrete\_Data} and create a data structure for it). Now - let's practice to achieve:

\subsection*{Class Statistics}

Compute:
\begin{itemize}
    \item Total number of samples
    \item Number in NSC
    \item Number in HSC
    \item the mean and standard deviation of the NSC concrete;
    \item the mean and standard deviation of the HSC concrete;
\end{itemize}

\begin{alert-box}{Python Programming Moment}
\begin{itemize}
    \item What Python packages are loaded and needed? 
    \item Reading the code, does it make sense to you?
\end{itemize}
\end{alert-box}


\subsection*{Histogram Visualization}

For selected features (e.g., Cement, Water, Age):

\begin{itemize}
    \item Plot histogram of each feature
    \item Overlay NSC and HSC distributions
\end{itemize}

\begin{alert-box}{\textbf{Questions to consider:}}
\begin{itemize}
    \item Do higher cement contents correlate with HSC?
    \item Does water content show inverse relationship?
    \item Does age significantly affect classification after standardization?
\end{itemize}
\end{alert-box}

\subsection*{Scatter Plots}

Plot two-feature combinations:

\begin{itemize}
    \item Cement vs Water
    \item Cement vs Age
    \item Water vs Superplasticizer
\end{itemize}

Color points by class label.

\begin{alert-box}{\textbf{Discussion:}}
\begin{itemize}
    \item Does a linear separation appear plausible?
    \item Are clusters evident?
\end{itemize}
\end{alert-box}

\section*{ML Production}

\subsection*{Train-Test Split}
We simulate the deployment of an AI-assisted mix screening tool.

\begin{itemize}
    \item 70\% training data
    \item 30\% testing data
    \item Use \textit{random shuffling}
\end{itemize}

\textbf{Engineering analogy:}

We do not certify safety using the same specimens used to calibrate predictive relationships.

\subsection*{Model Selection}
There are formal methods for model selection; at this moment, we will use a simple enumeration approach: use all the models we have learned so far. 

\subsubsection*{Model 1: Logistic Regression}

Logistic regression models probability as:

\[
P(y=1 \mid \mathbf{x}) = \sigma(\mathbf{w}^T \mathbf{x} + b)
\]

where:

\[
\sigma(z) = \frac{1}{1 + e^{-z}}
\]

\paragraph{Tasks}

\begin{enumerate}
    \item Train logistic regression on training set
    \item Compute:
    \begin{itemize}
        \item Training accuracy
        \item Testing accuracy
    \end{itemize}
\end{enumerate}

\subsubsection*{Model 2: Perceptron}

The perceptron performs linear classification using:

\[
y = \text{sign}(\mathbf{w}^T \mathbf{x} + b)
\]

\paragraph{Tasks}

\begin{enumerate}
    \item Train perceptron classifier
    \item Evaluate training and test accuracy
\end{enumerate}

\subsection*{Discussion}

\begin{itemize}
    \item Difference between logistic regression and perceptron?
    \item What loss functions are implicitly used?
\end{itemize}

\subsection*{Model Evaluation Metrics}

\subsection*{Accuracy}

\[
\text{Accuracy} = \frac{TP + TN}{TP + TN + FP + FN}
\]

\subsection*{Confusion Matrix}

\begin{center}
\begin{tabular}{c|cc}
 & Pred NSC & Pred HSC \\
\hline
Actual NSC & TP & FN \\
Actual HSC & FP & TN
\end{tabular}
\end{center}

\subsection*{Precision}

\[
\text{Precision} = \frac{TP}{TP + FP}
\]

\subsection*{Recall}

\[
\text{Recall} = \frac{TP}{TP + FN}
\]

\subsection*{Engineering Interpretation}

\begin{itemize}
    \item False Positive: HSC classified as NSC
    \item False Negative: NSC classified as HSC
\end{itemize}

Do we have a case where we need to determine which error is more critical?

\section*{Decision Boundary Visualization}

Choose two selected/critical features (e.g., Cement and Water); and choose two `seemingly' less critical features, do the same as follows:

\begin{enumerate}
    \item Plot decision boundary
    \item Overlay data points
\end{enumerate}

Do you observe or confirm your empirical knowledge (or engineering sense)?

\subsection*{Discussion}

\begin{itemize}
    \item Is boundary linear?
    \item Does it match engineering intuition?
\end{itemize}

\section*{Reflection Questions}

\begin{enumerate}
    \item How sensitive is classification to the 40 MPa threshold?
    \item Does standardization using ACI equation affect separability?
    \item Would adding regularization improve performance?
    \item How would you validate this model before field deployment?
\end{enumerate}

\section*{Summary}

In this tutorial we:

\begin{itemize}
    \item Standardized compressive strength to 28 days using ACI 209R-92
    \item Converted regression data to binary engineering classification
    \item Performed exploratory visualization
    \item Trained logistic regression and perceptron models
    \item Evaluated performance using accuracy, confusion matrix, precision, and recall
    \item Visualized decision boundaries
\end{itemize}

In the next, we will introduce:

\begin{itemize}
    \item Multi-layer Perceptrons (MLP)
    \item Linear SVM
    \item Kernel SVM
    \item Model selection and cross-validation; and advanced visual analytics.
\end{itemize}


\end{document}
